\section{Mathematical Formalization}


\subsection{Token Velocity Model}

Token velocity measures how often $\KEEPER$ changes hands:

\begin{equation}
V = \frac{T}{M}
\end{equation}

where:
\begin{itemize}
\item $V$: Velocity (transactions per token per year)
\item $T$: Total transaction volume (annual)
\item $M$: Money supply (circulating $\KEEPER$)
\end{itemize}

\textbf{Zoo Velocity Analysis:}

Transaction volume sources:
\begin{align}
T_{\text{inference}} &= N_{\text{daily}} \times F_{\text{avg}} \times 365 \\
T_{\text{experience}} &= U_{\text{exp}} \times P_{\text{exp}} \\
T_{\text{bridge}} &= B_{\text{volume}} \times 0.001 \\
T_{\text{total}} &= T_{\text{inference}} + T_{\text{experience}} + T_{\text{bridge}}
\end{align}

Base case (2025):
\begin{align}
T_{\text{inference}} &= 5M \times 0.005 \times 365 = 9.1B \text{ } \KEEPER \\
T_{\text{experience}} &= 100K \times 500K = 50B \text{ } \KEEPER \\
T_{\text{bridge}} &= 200B \times 0.001 = 0.2B \text{ } \KEEPER \\
T_{\text{total}} &= 59.3B \text{ } \KEEPER
\end{align}

Circulating supply (excluding staked):
\begin{equation}
M_{\text{circ}} = 1T \times (1 - s) = 1T \times 0.5 = 500B \text{ } \KEEPER
\end{equation}

(Assuming 50\% staking rate $s = 0.5$)

Velocity:
\begin{equation}
V = \frac{59.3B}{500B} = 0.119 \text{ (11.9\% annual turnover)}
\end{equation}

\textbf{Comparison to Other Cryptos:}
\begin{itemize}
\item Bitcoin: $V \approx 1.1$ (110\% annual turnover - high speculation)
\item Ethereum: $V \approx 5.2$ (520\% - DeFi usage drives velocity)
\item Stablecoins: $V \approx 12-20$ (high payment velocity)
\item Zoo: $V \approx 0.12$ (\textbf{very low - indicates hodling/staking})
\end{itemize}

Low velocity suggests $\KEEPER$ is \textbf{stored value} (staking, long-term holding) rather than transactional currency. This is ideal for governance token economics.

\subsection{Staking Equilibrium}

Optimal staking rate balances security (more staking) and liquidity (less staking):

\begin{equation}
s^* = \arg\max_{s} \left[ U_{\text{security}}(s) + U_{\text{liquidity}}(1-s) \right]
\end{equation}

where:
\begin{align}
U_{\text{security}}(s) &= \log(1 + s) \quad \text{(diminishing returns to security)} \\
U_{\text{liquidity}}(1-s) &= \log(1 + (1-s)) \quad \text{(diminishing returns to liquidity)}
\end{align}

Taking first-order condition:
\begin{equation}
\frac{dU}{ds} = \frac{1}{1+s} - \frac{1}{1+(1-s)} = 0
\end{equation}

Solving:
\begin{equation}
1 + s = 1 + (1 - s) \implies s^* = 0.5 \text{ (50\% optimal staking rate)}
\end{equation}

Zoo's design targets 40-60\% staking via APY adjustments:
\begin{itemize}
\item If $s < 0.4$: APY increases (more rewards per staked token) → incentivizes staking
\item If $s > 0.6$: APY decreases (rewards spread across more stakers) → incentivizes unstaking
\end{itemize}

\subsection{Experience Pricing Optimization}

Experience creators optimize price to maximize revenue:

\begin{equation}
\max_{p} R(p) = p \times D(p)
\end{equation}

where:
\begin{itemize}
\item $R(p)$: Revenue at price $p$
\item $D(p)$: Demand (usage count) as function of price
\end{itemize}

Assume demand follows power-law:
\begin{equation}
D(p) = D_0 \times p^{-\beta}
\end{equation}

where $\beta > 1$ is price elasticity.

Revenue:
\begin{equation}
R(p) = D_0 \times p^{1-\beta}
\end{equation}

Optimal price:
\begin{equation}
\frac{dR}{dp} = D_0 \times (1-\beta) \times p^{-\beta} = 0
\end{equation}

Since $(1-\beta) < 0$ for $\beta > 1$, revenue is \textbf{monotonically decreasing} in price. This suggests optimal strategy is \textbf{lowest viable price} to maximize volume.

However, this ignores quality signaling. Higher prices can signal higher quality:

\begin{equation}
D(p, q) = D_0 \times p^{-\beta} \times q^{\gamma}
\end{equation}

where $q$ = quality (effectiveness score), $\gamma > 0$.

Optimal price becomes:
\begin{equation}
p^* = \frac{\gamma}{\beta - 1} \times \frac{q}{1}
\end{equation}

Higher quality experiences can sustain higher prices proportional to $q^{\gamma}$.

